\section{Introduction}\label{sec:intro}

Decarbonizing the power system depends on demand that can move. Commercial buildings can cut or shift load when the grid is stressed, and wholesale demand-response (DR) programs pay them to do so \citep{siano2014,satchwell2021geb}. What is paid for must be measured, yet a reduction cannot be metered. The operator pays for the load a building shed: the load that would have occurred without the event minus the load actually metered (Section~\ref{sec:methods:metrics} gives the signed form). The first term needs an estimate, the customer baseline load (CBL), and the calculation that turns it into a payment is the settlement. A baseline comes from an algorithm over a defined data stream, so the hours that algorithm reads decide how it can be wrong. In computing terms, it is a counterfactual estimator whose output sets a payment and whose ground truth is unobserved under a real event. It divides money among parties and is not a forecast that serves control \citep{li2026review}. Its signed bias sets the direction of any systematic transfer between the parties, and its absolute error sets the size of the transfer \citep{wijaya2014}. Evaluating baselines is the measurement layer of building demand flexibility, which places it within the scope of this journal: computing for sustainability, and the modeling and evaluation of sustainable systems.

This journal has published related work from the computing side. Zhang et al.\ \citep{zhang2024suscom} coordinate data centers and a load aggregator against a target load reduction. Their contract is a shortfall penalty on simulated counterfactuals, so it needs no baseline estimate and represents no baseline error. Abdelghani et al.\ \citep{abdelghani2024suscom} predict building occupancy from Wi-Fi logs, and Wang et al.\ \citep{wang2023suscom} apply a machine-learning optimization model to reduce electricity loads in buildings. Neither settles a payment. Kannan et al.\ \citep{kannan2025suscom} forecast load for smart-grid DR, and He et al.\ \citep{he2025suscom} build a digital twin of smart-grid operations with simulated gains. None of these studies evaluates, on real meters and with intervals, the counterfactual on which a settlement pays.

The DR literature has measured baselines for over a decade. Coughlin et al.\ compared seven baselines on proxy event days at California non-residential sites \citep{coughlin2009}. Granderson et al.\ found that baselines under-predicted peak windows for most of 453 commercial buildings, and that same-day adjustment reduced the bias \citep{granderson2021}. Control groups and synthetic controls were tested on residential data \citep{hatton2016,wang2018spm,lee2019,sievers2026} and proposed as a settlement option in California \citep{caiso_bawg2017}. Pre-cooling has been reported to inflate adjusted baselines \citep{coughlin2009,braithwait2010,li2026review}. Frequent dispatch pushes the lookback further back \citep{enernoc2009naesb,dsa2023elrp}, an effect Granderson et al.\ examined by sliding the lookback on non-DR data \citep{granderson2023}. Todd et al.\ show why the truth is hard to find: in a residential program, spillover made a common baseline understate event reductions \citep{todd2019}.

Rule accuracy on proxy event days and the value of same-day adjustment are known. Field reports, policy documents and simulations show qualitatively that a pre-event load increase inflates an adjusted baseline and that frequent dispatch pushes the lookback further back \citep{gabaldon2021,afentoulis2025,wang2022manipulation,wang2026physics}. What is not known is how large an increase, and how many dispatches, each information set tolerates on real campus and commercial meters. A 2026 review names this gap and asks researchers to ``evaluate the manipulation-resistance of various estimation methods'' \citep{li2026review}. Prior evidence has three limitations: real events do not reveal the truth \citep{todd2019}, earlier commercial meter data were proprietary or held by the authors \citep{granderson2021,mathieu2011cdc}, and the comparison-group studies use residential samples. Known-truth injection into open meter data can quantify the sensitivity of each settlement formula to specified perturbations, although it observes neither real dispatch nor participant behavior.

We build a measurement instrument, not a new estimator: a reproducible pipeline that maps baseline information sets to settlement errors under controlled load transformations. An information set is the data stream a baseline reads. It is the building's own history (I-hist), the metered hours of the event day before the window (I-morn), or the same-day load of other buildings (I-donor). The pipeline injects known changes into hourly meters of the Building Data Genome Project 2 (BDG2) \citep{miller2020bdg2}, so that we know the injected change, and hence the settlement error, by construction. It covers {{sample.n_buildings|int}} buildings at {{sample.n_sites|int}} campus-heavy sites, uses hot weekdays as proxy events, and follows a pre-specified and hash-frozen analysis plan with a 2016 development year and a 2017 held-out year at the same sites (Fig.~\ref{fig:design}). We ask three questions.

\begin{itemize}
\item \textbf{RQ1.} How do hot-day error and bias depend on the information set a baseline reads (hypotheses H1 to H3)?
\item \textbf{RQ2.} How large an injected four-hour pre-event load increase does the day-of-adjusted rule of PJM Interconnection (PJM-adj) tolerate before a baseline that does not read those hours becomes more accurate (H4)?
\item \textbf{RQ3.} How does the absolute estimation error of curtailment grow under repeated dispatch (H5)?
\end{itemize}

The study makes four contributions, each resting on stress tests of settlement formulas on injected load and not on observed participant behavior.

\begin{enumerate}
\item[(i)] A measurement of the break-even pre-event increase, the injected four-hour increase at which PJM-adj loses its hot-day accuracy advantage over a ridge synthetic control (SCM-ridge) that does not read those hours. The four-hour statistic is the median of the site break-evens, {{H4.median_site_pcross|pct1}}, with a site-level sign test against the pre-set ten percent threshold. Because the shaped hours cover all three PJM adjustment hours, it is a worst-case placement bound. The two-hour result is building-pooled, with no site-level test, and its interval includes the threshold.
\item[(ii)] A trade-off across the two stress tests. In the injected repeated-dispatch stress test, PJM-adj stayed more accurate than SCM-ridge at each tested dispatch count in both event windows. In the confirmatory 14:00 to 18:00 window its absolute error also grew more slowly (slope difference {{H5.eta_diff|.4f}}), but it still grew; in the pre-specified secondary 16:00 to 20:00 window the slope contrast was inconclusive. Sensitivity to shaping follows the overlap of read and shapeable hours, and sensitivity to dispatch follows how far dispatch pushes back the data a baseline reads.
\item[(iii)] A site-level equivalence result. Under the signed-rank location (symmetry) assumption, the adjusted comparison group (CG-adj), which is PJM-adj with a donor ratio, was equivalent to PJM-adj at the site median within the pre-set ten percent margin. Yet {{derived.H3.n_sites_outside_margin|int}} of {{sample.n_sites|int}} sites lay outside it, and the conditional bootstrap interval crossed the lower margin.
\item[(iv)] An open pipeline on public meter data, with a pre-specified, hash-frozen analysis plan and exact site-level inference. It is measurement infrastructure, not a new algorithm.
\end{enumerate}

Contributions (i) and (ii) bear on the open problem of testing the manipulation-resistance of estimators, and (ii) and (iv) on the review's concern that frequent dispatch leaves few non-DR days for baselines \citep{li2026review}. Two known results are replications, not contributions: day-of adjustment reduces hot-day error, and the hot-day bias of unadjusted PJM is more negative than its random-weekday bias.

Section~\ref{sec:background} reviews rules and prior evidence, Section~\ref{sec:data} the data and design, and Section~\ref{sec:methods} the baselines, stress tests and inference. Section~\ref{sec:results} gives results, Section~\ref{sec:discussion} their meaning and limits, and Section~\ref{sec:conclusion} conclusions.
